\chapter{Problem Definition, Course Alignment, and Scope}
\label{ch:scope}

\section{Clinical and computational motivation}

Alzheimer's disease and frontotemporal dementia can present overlapping
cognitive and behavioral symptoms, while their electrophysiological signatures
may vary across people, time, frequency, and scalp location. Resting-state
\EEG{} is non-invasive and comparatively accessible, but its high temporal
resolution creates a modeling problem: a single recording is not naturally a
flat table. It has participant, channel, time, frequency, and potentially
connectivity modes. The course premise that physiological data can be organized
as \emph{channels $\times$ time $\times$ frequency $\times$ subject} therefore
maps directly to this dataset \citep{padhy2026course}.

The project asks a concrete question:

\begin{quote}
Can carefully tensorized, low-rank representations of artifact-cleaned,
eyes-closed resting \EEG{} support reproducible subject-level classification of
\AD{}, \FTD{}, and \CN{} while preserving multi-way structure and avoiding
window-level leakage?
\end{quote}

The answer is investigated comparatively rather than assumed. Matrix baselines,
non-negative CP/Tucker, Tensor Train, waveform-delay Tucker/TT, a supervised TT
variant, spectral controls, and several classifier families are all retained in
the record. This makes the project as much an evaluation of representation
choices as a search for a strong classifier.

\section{Course learning outcomes addressed}

The completed work operationalizes the course concepts in four ways:

\begin{enumerate}
  \item \textbf{Tensorization as a modeling decision.} The same \EEG{} source is
  represented as a subject-channel-frequency-time tensor, a signed
  channel-frequency-local-segment tensor, and a channel-lag-position delay
  tensor. These are not cosmetic reshapes; each encodes a different invariance.
  \item \textbf{Low-rank structure.} CP, Tucker/HOSVD, and TT decompositions
  replace large mode products with compact factors or reusable bases, following
  the rank-one and multilinear ideas emphasized in Module 2
  \citep{padhy2026lowrank,kolda2009tensor,oseledets2011tt}.
  \item \textbf{Multimodal interactions.} Channel, frequency, local time, lag,
  and participant modes are kept explicit long enough for decomposition to
  exploit their interactions.
  \item \textbf{Responsible model comparison.} The project demonstrates why
  tensor structure, optimization convergence, and prediction accuracy are
  distinct questions. Strong reconstruction or a converged optimizer does not
  guarantee strong held-out clinical classification.
\end{enumerate}

\section{Proposal-to-implementation traceability}

The one-page proposal described an ambitious graph-regularized tensor biomarker
system. Table~\ref{tab:traceability} distinguishes completed work from planned
but unimplemented extensions. This prevents proposal language from being
mistaken for results.

\begin{longtable}{p{0.27\textwidth}p{0.31\textwidth}p{0.34\textwidth}}
\caption{Traceability from the formal proposal to the implemented workspace.}
\label{tab:traceability}\\
\toprule
\textbf{Proposal element} & \textbf{Implemented counterpart} & \textbf{Status and interpretation}\\
\midrule
\endfirsthead
\toprule
\textbf{Proposal element} & \textbf{Implemented counterpart} & \textbf{Status and interpretation}\\
\midrule
\endhead
Subject $\times$ channel $\times$ frequency $\times$ time tensor &
V1 tensor of shape $69\times19\times59\times10$ and V2 tensors of shape
$19\times59\times3$ per window & Completed, with Welch/spectrogram PSD rather
than Morlet wavelets.\\
Dynamic phase-locking value connectivity & No dPLV feature branch in the saved
experiments & Not implemented; no connectivity-result claim is made.\\
Graph-regularized non-negative CP/Tucker & Preliminary graph NCP in V1; corrected
fixed anatomical pre-smoothing before TT in V2 & Partially implemented. The old
graph-NCP objective had a scale-consistency caveat and remains exploratory.\\
CORCONDIA and uniqueness analysis & Rank sweeps, reconstruction diagnostics,
orthogonality tests, and numerical TT checks & CORCONDIA itself was not
implemented.\\
LOSO-CV & Repeated five-fold subject CV in V1 and nested repeated five-fold
subject CV in V2--V4 & Replaced with repeated nested subject CV to make bounded
model/rank selection feasible.\\
PSD and Riemannian baselines & PSD/spectral controls in every relevant stage &
PSD completed; Riemannian covariance baseline not implemented.\\
MMSE correlation and biomarker topoplots & MMSE retained in metadata but excluded
from predictors; no validated factor-MMSE analysis & Not implemented; this
report does not imply clinical severity biomarkers.\\
Three-way AD/FTD/CN classifier & Logistic regression, SVM, histogram boosting,
random forest, XGBoost, and supervised TT & Completed as subject-level
comparative experiments.\\
\bottomrule
\end{longtable}

\section{Team and proposal-assigned modules}

The formal proposal assigned the following course modules. Wording here is
normalized for the report, while the names, roll numbers, and allocation follow
the supplied proposal.

\begin{table}[htbp]
\centering
\small
\begin{tabularx}{\textwidth}{p{0.25\textwidth}p{0.14\textwidth}X}
\toprule
\textbf{Member} & \textbf{Roll number} & \textbf{Proposal-assigned technical module}\\
\midrule
Neeruganti Jeevan Sai & 24bds046 & Multi-way tensor construction, time-frequency representation, synchrony-array design, and normalization.\\
Kothapalli Tharun Chowdary & 24bds032 & Optimization and regularization, graph-aware CP/Tucker design, and electrode adjacency construction.\\
Vatipalli Abhiram & 24bds087 & Rank selection and stability, convergence analysis, hyperparameter exploration, and residual profiling.\\
Palyam Maruthi & 24bcs098 & Biomarker mapping and statistics, group-level comparison, topographic projection, and MMSE association planning.\\
Shaik Moiz & 24bcs133 & Classification and benchmarking with SVM/logistic models, metric analysis, and baseline comparison.\\
\bottomrule
\end{tabularx}
\caption{Team list and planned module allocation from the formal proposal.}
\label{tab:team}
\end{table}

\section{Research questions and success criteria}

The completed pipeline addresses five research questions:

\begin{enumerate}
  \item Which diagnosis limits balanced participant and duration coverage?
  \item Can a reproducible primary cohort be built without discarding excluded
  windows or synthetically oversampling \FTD{}?
  \item Which tensorization preserves enough discriminative structure for
  subject-level prediction?
  \item Do TT/Tucker representations improve over spectral controls when every
  learned step is confined to training people?
  \item How much does bounded classifier and optimization tuning improve the
  complete procedure, and which diagnosis remains difficult?
\end{enumerate}

Success is not defined as reaching an arbitrary 65--70\% threshold. A credible
course result must instead be reproducible, leakage-aware, explicit about the
primary endpoint, complete about negative findings, and honest about the lack of
external validation.

